\subsection{The literal integral localization quotient}
\label{sec:ghost-strict-localization}

\begin{theorem}[Integral localization with the coefficient maps retained]
\label{thm:ghost-strict-localization}
For every prime \(\ell\) and every integer \(r\ge2\), put
\(T=\mathbb Z_\ell(r)\), \(S=\{\ell,\infty\}\), and
\[
 R_\ell=\{(z_q)_{q\ne\ell}\in\prod_{q\ne\ell}\mathbb Z_\ell:
                 z_q\longrightarrow0\},\qquad
 (d_rz)_q=(q^{r-1}-1)z_q.
\]
Fix compatible Tate coordinates at the finite residue fields.
There is a natural, in those coordinates, exact sequence of
abelian groups
\[
 0\longrightarrow H^2_{\rm cts}(G_{\mathbb Q,S},T)
 \xrightarrow{\rm inf}H^2_{\rm cts}(G_{\mathbb Q},T)
 \xrightarrow{\Psi_{\ell,r}}R_\ell/d_rR_\ell
 \longrightarrow0.
\]
The quotient uses the actual image \(d_rR_\ell\).
In particular this proves the exactness assertion used in
SECTIONS A.65.8. For \(r=3\), inflation has zero source,
so \(\Psi_{\ell,3}\) is an isomorphism for every \(\ell\),
as asserted abstractly in SECTIONS A.46.2.
\end{theorem}
\begin{proof}
We specify the morphism before taking any inverse limit.
Write \(T_n=T/\ell^nT\) and
\(X=\operatorname{Spec}\mathbb Z[1/\ell]\). For a finite
set \(F\) of closed points of \(X\), the sequence of
cochains with supports, or its associated distinguished
triangle in the derived category of abelian groups, is
\[
 R\Gamma(X,T_n)\longrightarrow
 R\Gamma(X\setminus F,T_n)\longrightarrow
 \bigoplus_{q\in F}R\Gamma(\mathbb F_q,T_n(-1))[-1]
 \longrightarrow R\Gamma(X,T_n)[1].
\]
Purity is precisely the calculation
\(R^2i^!\mu_{\ell^n}=\mathbb Z/\ell^n\), with all
other terms zero, in \cite[II.1.7(b)]{GhostMilneADT};
tensoring with the locally free twist \(\mu_{\ell^n}^{\otimes(r-1)}\)
gives the displayed twist and shift. The identifications
are natural under coefficient reduction and under enlarging
\(F\). Filtered direct limits are exact in abelian groups.
Taking their limit and using the group-to-scheme comparison
\cite[II.2.9]{GhostMilneADT} yields
\[
 \begin{gathered}
 R\Gamma(G_{\mathbb Q,S},T_n)\longrightarrow
 R\Gamma(G_{\mathbb Q},T_n)\longrightarrow L_n[-1]
 \longrightarrow R\Gamma(G_{\mathbb Q,S},T_n)[1],\\
 L_n=\bigoplus_{q\ne\ell}R\Gamma(\mathbb F_q,T_n(-1)).
 \end{gathered}
\]
One can take the middle term to be the usual continuous
Galois cochain complex: each finite-valued cochain factors
through a finite Galois quotient, hence is defined over
some \(X\setminus F\). Thus this triangle concerns the
absolute Galois group appearing in the statement.

Arithmetic Frobenius topologically generates
\(G_{\mathbb F_q}=\widehat{\mathbb Z}\) and acts on
\(T_n(-1)\) by \(q^{r-1}\). The procyclic resolution
therefore computes this cohomology by the two-term complex
\[
 [\,\mathbb Z/\ell^n
          \xrightarrow{q^{r-1}-1}\mathbb Z/\ell^n\,]
                     \quad\hbox{in degrees zero and one}.
\]
For clarity, the prime-to-\(\ell\) subgroup of
\(\widehat{\mathbb Z}\) acts on this rank-one Tate
module through a finite group of order prime to \(\ell\).
When its action is nontrivial, \(q^{r-1}-1\) is a unit
in \(\mathbb Z_\ell\). Its invariant module is zero,
and the full displayed two-term complex is compatibly
contractible, by multiplication by that unit's inverse.
When the prime-to-\(\ell\) action is trivial, taking
its invariants preserves the whole module. The remaining
pro-\(\ell\) cyclic resolution is the sequence
\(0\to\mathbb Z_\ell[[u]]\xrightarrow{u}
\mathbb Z_\ell[[u]]\to\mathbb Z_\ell\to0\);
the invariant functor just used is exact by compatible
normalized averaging. Both cases are functorial in the
coefficient module, so the two-term models commute with
reduction modulo \(\ell^n\). Consequently \(L_n\) is
represented compatibly by
\[
 [\,\bigoplus_{q\ne\ell}\mathbb Z/\ell^n
       \xrightarrow{d_r}\bigoplus_{q\ne\ell}\mathbb Z/\ell^n\,].
\]
Both terms have surjective coefficient reductions. Their
inverse limit is exactly \(R_\ell\): a compatible family
has finite support modulo each \(\ell^n\), which is the
displayed convergence condition, and conversely that
condition gives such a family. The derived inverse limit
of \(L_n\) is therefore the literal two-term complex
\([R_\ell\xrightarrow{d_r}R_\ell]\).

Here derived inverse limit can be constructed without any
topological closure operation: for a tower of complexes
\(C_n\), it is the mapping fibre of one minus the shift on
\(\prod_n C_n\). This construction carries the compatible
triangles above to a triangle. For the global Galois
cochain complexes, coefficient reduction is surjective in
each degree, because a finite-valued continuous function
lifts by lifting each of its finitely many values.
Their derived inverse limits are thus the usual continuous
cochain complexes with coefficients in \(T\). The same
holds for \(G_{\mathbb Q,S}\). Taking the degree-two
part gives
\[
 \ker d_r\longrightarrow H^2_{\rm cts}(G_{\mathbb Q,S},T)
 \longrightarrow H^2_{\rm cts}(G_{\mathbb Q},T)
 \xrightarrow{\Psi_{\ell,r}}\operatorname{coker}d_r
 \longrightarrow H^3_{\rm cts}(G_{\mathbb Q,S},T)
 \longrightarrow H^3_{\rm cts}(G_{\mathbb Q},T).
\]
This is the promised definition of the residue morphism
\(\Psi_{\ell,r}\), induced by the compatible supports
morphisms. Since each nonzero integer \(q^{r-1}-1\)
acts injectively on \(\mathbb Z_\ell\), its preceding
kernel is zero.

The last displayed inflation is injective. Indeed,
finite coefficient \(H^2(G_{\mathbb Q,S},T_n)\) is finite
by finite-place global duality and local finiteness
\cite[I.2.1, I.4.10]{GhostMilneADT}; its derived inverse
limit is zero. Finite coefficient \(H^3\) is identified
by real restriction with \(H^3(G_{\mathbb R},T_n)\)
by I.4.10(c). At the real group the degree-two groups
are also finite. Thus real restriction induces an
isomorphism on integral degree three from
\(G_{\mathbb Q,S}\). Its factorization through
\(G_{\mathbb Q}\) proves the required injectivity,
including \(\ell=2\) and odd \(r\), where the real
group may be nonzero. Exactness now gives the asserted
short exact sequence.

Finally \(K_4(\mathbb Z)=0\)
\cite[VI.10.1.1]{GhostWeibelKbook}. The odd-primary
comparison VI.8.2 and the two-primary comparison
VI.9.11, case \(n=4\), give
\(H^2(G_{\mathbb Q,\{\ell,\infty\}},\mathbb Z_\ell(3))=0\)
for every \(\ell\). The group-to-scheme comparison and
the vanishing of the finite degree-one derived limit are
the same as above. This proves the final isomorphism.
\end{proof}

The theorem does not assign an arbitrary representative
of the local finite classes to an integral obstruction.
Its morphism is the inverse limit of the actual supports
morphisms. In particular, to use the formulas A.46.5 or
A.65.12, one must still identify their chosen tame-relator
cochains with this morphism. That explicit identification
is not supplied merely by the quotient's abstract
description.
